\section*{الفصل \ref{ch:g11:fossil-fuel-geology} --- جيولوجيا الفحم والنفط والغاز}

\begin{solution}{exo:g11:fossil-fuel-geology:1}
تستقر العوالق الميتة على قاع بحر فقير بالأكسجين؛ تصير المادة العضوية
كيروجينًا في وحل؛ الدفن تحت رواسب أحدث؛ \omterm{def:g11:fossil-fuel-geology:oil-window}{النضج} في
\omterm{def:g11:fossil-fuel-geology:oil-window}{نافذة النفط}؛ \omterm{def:g11:fossil-fuel-geology:migration}{الهجرة}؛ يتجمع النفط تحت \omterm{def:g11:fossil-fuel-geology:reservoir}{صخر حاجز}.
\end{solution}

\begin{solution}{exo:g11:fossil-fuel-geology:2}
(1) \omterm{def:g11:fossil-fuel-geology:source-rock}{صخر مصدر}. (2) \omterm{def:g11:fossil-fuel-geology:reservoir}{صخر خازن}. (3) \omterm{def:g11:fossil-fuel-geology:reservoir}{صخر حاجز}. (4) \omterm{def:g11:fossil-fuel-geology:reservoir}{صخر حاجز}. (5) \omterm{def:g11:fossil-fuel-geology:reservoir}{صخر خازن}.
\end{solution}

\begin{solution}{exo:g11:fossil-fuel-geology:3}
الأعلى: $(60 - 10)/25 = \qty{2.0}{km}$؛ والقاعدة: $(120 - 10)/25 = \qty{4.4}{km}$.
وسمك النافذة $60/25 = \qty{2.4}{km}$.
\end{solution}

\begin{solution}{exo:g11:fossil-fuel-geology:4}
(1) الخث، فالليغنيت، \omterm{def:g11:fossil-fuel-geology:coal}{فالفحم} البيتوميني، فالأنثراسيت. (2) الأنثراسيت، من 86 إلى
\qty{97}{\percent} من الكربون. (3) الأنثراسيت، فيما وراء نحو \qty{5}{km}
و\qty{150}{\degreeCelsius}.
\end{solution}

\begin{solution}{exo:g11:fossil-fuel-geology:5}
(1) تتمايز الموائع بحسب كثافتها: الغاز أخفها، ثم النفط،
ثم الماء. (2) يصعد النفط والغاز لأنهما أخف من الماء؛ ولا
يوقفهما إلا صخر لا يستطيعان عبوره. الحجر الطيني (أو السجيل) والملح.
\end{solution}

\begin{solution}{exo:g11:fossil-fuel-geology:6}
(1) A: $(60 - 15)/20 = \qty{2.25}{km}$ و$(120 - 15)/20 = \qty{5.25}{km}$.
B: $45/40 \approx \qty{1.13}{km}$ و$105/40 \approx \qty{2.63}{km}$.
(2) على عمق \qty{2}{km}: في A، $15 + 40 = \qty{55}{\degreeCelsius}$، فهو لم ينضج
بعد، ولا نفط؛ وفي B، $15 + 80 = \qty{95}{\degreeCelsius}$، في \omterm{def:g11:fossil-fuel-geology:oil-window}{نافذة
النفط}، فهو يصنع النفط. (3) في الحوض A: لا يتكوّن النفط فيه إلا تحت
\qty{2.25}{km}.
\end{solution}

\begin{solution}{exo:g11:fossil-fuel-geology:7}
(1) $11 \times 3 = \qty{33}{m}$ من الخث. (2) $\qty{33000}{mm}$ بمعدل
\qty{1}{mm} في السنة: نحو \num{33000} سنة. (3) لا يبقى الخث إلا حيث
يكون تحت ماء راكد، بعيدًا عن متناول الأكسجين؛ وكلما تراكم، وجب أن
تهبط الأرض بالوتيرة نفسها كي يبقى سطحه مشبعًا بالماء.
\end{solution}

\begin{solution}{exo:g11:fossil-fuel-geology:8}
(1) $20 + 30 \times 1.2 = \qty{56}{\degreeCelsius}$؛
$20 + 30 \times 2.3 = \qty{89}{\degreeCelsius}$؛
$20 + 30 \times 4.0 = \qty{140}{\degreeCelsius}$.
(2) قبل 120 مليون سنة، غير ناضج؛ وقبل 95 مليون سنة، في \omterm{def:g11:fossil-fuel-geology:oil-window}{نافذة
النفط}؛ وقبل 20 مليون سنة، في \omterm{def:g11:fossil-fuel-geology:oil-window}{نافذة الغاز}.
\end{solution}

\begin{solution}{exo:g11:fossil-fuel-geology:9}
(1) $2580 - 2400 = \qty{180}{m}$. (2) يملأ الغاز
$2450 - 2400 = \qty{50}{m}$، والنفط $2580 - 2450 = \qty{130}{m}$. (3) ينسكبان
خارجها عند نقطة الانسكاب، تحت أعمق خط تسوية منغلق، ويواصلان \omterm{def:g11:fossil-fuel-geology:migration}{الهجرة}
صعودًا، إلى \omterm{def:g11:fossil-fuel-geology:trap}{مصيدة} أخرى أو إلى السطح.
\end{solution}

\begin{solution}{exo:g11:fossil-fuel-geology:10}
(1) أعلى النافذة عند \qty{3}{km}: $(60 - 20)/3 \approx
\qty{13}{\degreeCelsius/km}$؛ والقاعدة: $(120 - 20)/3 \approx
\qty{33}{\degreeCelsius/km}$. (2) على عمق \qty{3}{km} تكون درجة الحرارة من
$20 + 75 = \qty{95}{\degreeCelsius}$ إلى $20 + 102 = \qty{122}{\degreeCelsius}$:
في \omterm{def:g11:fossil-fuel-geology:oil-window}{نافذة النفط}، أو عند قاعدتها مع أشد التدرجات انحدارًا. والقاعدة
معقولة: ففي معظم الأحواض يكون \omterm{def:g11:fossil-fuel-geology:source-rock}{الصخر المصدر} المدفون تحت \qty{3}{km} من الصخور قد
بدأ بصنع النفط.
\end{solution}

\begin{solution}{exo:g11:fossil-fuel-geology:11}
(1) \qty{30}{\percent}: ليغنيت؛ و\qty{70}{\percent}: \omterm{def:g11:fossil-fuel-geology:coal}{فحم} بيتوميني؛
و\qty{92}{\percent}: أنثراسيت. (2) الأنثراسيت، فيما وراء نحو
\qty{5}{km} و\qty{150}{\degreeCelsius}. (3) أسود زجاجي.
\end{solution}

\begin{solution}{exo:g11:fossil-fuel-geology:12}
(1) $1100/25 = \qty{44}{m}$ في كل مليون سنة. (2) $0.044 \times 30
\approx \qty{1.3}{\degreeCelsius}$ في كل مليون سنة. (3) تمتد النافذة على
\qty{60}{\degreeCelsius}: $60/1.32 \approx 45$ مليون سنة. (4) تباطأ الدفن
بعد 95 مليون سنة مضت (نحو 24، ثم 23
و\qty{20}{m} في كل مليون سنة): فسخن \omterm{def:g11:fossil-fuel-geology:source-rock}{الصخر المصدر} ببطء أكبر.
\end{solution}

\begin{solution}{exo:g11:fossil-fuel-geology:13}
(1) $(80 - 51)/(117 - 51) = 29/66 \approx 0.44$: نحو \qty{44}{\percent}.
(2) هاجر صعودًا قبل أن توجد فوقه أي \omterm{def:g11:fossil-fuel-geology:trap}{مصيدة}: فأفلت إلى
السطح، حيث تأكسد، أو تجمع في \omterm{def:g11:fossil-fuel-geology:trap}{مصيدة} أخرى. (3) قبل 40
مليون سنة كان \omterm{def:g11:fossil-fuel-geology:source-rock}{الصخر المصدر} قد غادر \omterm{def:g11:fossil-fuel-geology:oil-window}{نافذة النفط}: فلم يكن يصنع إلا الغاز،
ولم يكن في وسع الطية المحدبة أن تحتجز منه إلا الغاز.
\end{solution}

\begin{solution}{exo:g11:fossil-fuel-geology:14}
(1) الأعلى: $(60 - 20)/40 = \qty{1.0}{km}$؛ والقاعدة: $(120 - 20)/40 =
\qty{2.5}{km}$. (2) يقع \qty{1.0}{km} بين \qty{0.4}{km} (قبل 150 مليون
سنة) و\qty{1.2}{km} (قبل 120): $150 - 30 \times 0.6/0.8 = 127.5$
مليون سنة مضت. ويقع \qty{2.5}{km} بين \qty{2.3}{km} (قبل 95)
و\qty{3.0}{km} (قبل 66): $95 - 29 \times 0.2/0.7 \approx 86.7$ مليون سنة
مضت. (3) لا نفط: فقد غادر \omterm{def:g11:fossil-fuel-geology:source-rock}{الصخر المصدر} \omterm{def:g11:fossil-fuel-geology:oil-window}{نافذة النفط} قبل أن تتكوّن الطية
المحدبة، وكان نفطه قد هاجر بعيدًا. ويمكنها أن تحتجز الغاز المتكوّن بعد 80
مليون سنة مضت.
\end{solution}

\begin{solution}{exo:g11:fossil-fuel-geology:15}
(1) \qty{5}{km} على الأقل. (2) $15 + 30 \times 5 = \qty{165}{\degreeCelsius}$،
وهي فوق \qty{150}{\degreeCelsius}: فهي متفقة. (3) نُهض \qty{5}{km} على الأقل
من الصخور ثم تآكلت. و$\qty{5e6}{mm}$ في \num{299} مليون سنة على الأكثر
تعطي على الأقل $5 \times 10^6/2.99 \times 10^8
\approx \qty{0.017}{mm}$ في السنة في المتوسط.
\end{solution}

\begin{solution}{pb:g11:fossil-fuel-geology:1}
\textbf{1.} \omterm{def:g11:fossil-fuel-geology:source-rock}{الصخر المصدر}: السجيل الأسود؛ \omterm{def:g11:fossil-fuel-geology:reservoir}{والصخر الخازن}: الحجر الرملي المسامي؛
\omterm{def:g11:fossil-fuel-geology:reservoir}{والصخر الحاجز}: الحجر الطيني الذي فوقه.

\textbf{2.} عوالق استقرت على قاع بحر فقير بالأكسجين فأفلتت من
التحلل ودُفنت في الوحل.

\textbf{3.} \omterm{def:g11:fossil-fuel-geology:source-rock}{الصخر المصدر} في الجوراسي (فالعدد 165 يقع بين 201.4
و143.1). والحجر الرملي بدأ في الجوراسي وترسب معظمه في
الطباشيري المبكر.

\textbf{4.} المادة العضوية المبعثرة في صخر رسوبي، التي لا تذوب في
المذيبات العضوية. ولا تصنع النفط إلا إذا سُخّنت إلى نحو
\qty{60}{\degreeCelsius}، وهذا يتطلب الدفن.

\textbf{5.} $400 + 800 + 1100 + 700 + 700 + 300 - 500 = \qty{3500}{m}$.

\textbf{6.} \qty{0.4}{km}، \qty{1.2}{km}، \qty{2.3}{km}، \qty{3.0}{km}،
\qty{3.7}{km}، \qty{4.0}{km}؛ واليوم \qty{3.5}{km}.

\textbf{7.} $400/15 \approx 26.7$؛ $800/30 \approx 26.7$؛ $1100/25 = 44$؛
$700/29 \approx 24.1$؛ $700/31 \approx 22.6$؛ $300/15 = 20$ (\unit{m} في كل
مليون سنة).

\textbf{8.} قبل 20 مليون سنة، على عمق \qty{4.0}{km}.

\textbf{9.} نهض الحوض وتآكل \qty{500}{m} من الصخور من
أعلاه.

\textbf{10.} ترسبت قاعدة الحجر الرملي قبل 150 مليون سنة، فوق
\qty{400}{m} من السجيل والحجر الطيني: $3.5 - 0.4 = \qty{3.1}{km}$.

\textbf{11.} $T = 20 + 30 z$: \qty{32}{\degreeCelsius}، \num{56}، \num{89}،
\num{110}، \num{131}، \qty{140}{\degreeCelsius}؛ واليوم
\qty{125}{\degreeCelsius}.

\textbf{12.} $40/30 \approx \qty{1.33}{km}$ و$100/30 \approx
\qty{3.33}{km}$.

\textbf{13.} بين 120 مليون سنة مضت (\qty{56}{\degreeCelsius})
و95 مليون سنة مضت (\qty{89}{\degreeCelsius}).

\textbf{14.} ارتفعت درجة الحرارة 33 درجة في 25 مليون سنة؛ فاستغرق
ارتفاعها الدرجات الأربع من \num{56} إلى \qty{60}{\degreeCelsius}
$25 \times 4/33 \approx 3.0$ مليون سنة: أي نحو $120 - 3.0 = 117$ مليون سنة
مضت.

\textbf{15.} بين 66 مليون سنة مضت (\qty{110}{\degreeCelsius})
و35 (\qty{131}{\degreeCelsius}): $66 - 31 \times 10/21 \approx 51.2$ مليون
سنة مضت.

\textbf{16.} \qty{140}{\degreeCelsius}، قبل 20 مليون سنة. وعند
\qty{125}{\degreeCelsius} اليوم هو في \omterm{def:g11:fossil-fuel-geology:oil-window}{نافذة الغاز}: يصنع الميثان في
الغالب.

\textbf{17.} لم تكن فوقه \omterm{def:g11:fossil-fuel-geology:trap}{مصيدة}، فهاجر صعودًا وأفلت، إلى
السطح أو إلى \omterm{def:g11:fossil-fuel-geology:trap}{مصائد} أخرى.

\textbf{18.} $(80 - 51.2)/(117.0 - 51.2) \approx 0.44$.

\textbf{19.} الغاز في الأعلى، ثم النفط، ثم الماء.

\textbf{20.} دخل \omterm{def:g11:fossil-fuel-geology:source-rock}{الصخر المصدر} \omterm{def:g11:fossil-fuel-geology:oil-window}{نافذة النفط} قبل نحو \textbf{117
مليون سنة}.
\end{solution}
